August 3, 2026

Editor-in-Chief
IEEE Access
IEEE

SUBJECT: Submission of Manuscript P5 — Hardware Efficiency Modeling and Adaptive Thermal Throttling in Edge Sensing Nodes

Dear Editor-in-Chief and Associate Editors,

We are pleased to submit our original research manuscript titled "Hardware Efficiency Modeling and Adaptive Thermal Throttling in Edge Sensing Nodes" for publication consideration in IEEE Access.

1. SUBSYSTEM BACKGROUND AND MOTIVATION
Deploying multi-stream deep learning models (face recognition and pose estimation) on edge hardware devices (such as NVIDIA Jetson Orin Nano or Apple Silicon M2 Mac mini) under 24/7 continuous operational workloads introduces severe thermal management challenges. Sustained peak GPU utilization leads to thermal throttling, frame dropping, and hardware failure. This paper models and mitigates thermal stress in edge-native sensing platforms.

2. CORE TECHNICAL CONTRIBUTIONS AND NOVELTY
  - Adaptive Thermal Throttling Engine: Implements a hardware monitoring daemon (`PowerThread`) that continuously samples system temperatures via `psutil`, dynamically scaling video ingestion rates from 30 FPS to 15 FPS when temperatures cross 85°C.
  - 24-Hour Continuous Load Benchmark: Provides comprehensive empirical profiling of CPU/GPU thermal dissipation, power consumption, and memory stability under continuous multi-camera workloads.
  - Failure Prevention & Stability: Eliminates thermal-induced system crashes and frame buffer corruption, maintaining system availability without active liquid cooling.

3. SUITABILITY FOR IEEE ACCESS
This manuscript aligns with IEEE Access's scope in practical hardware engineering, edge computing performance optimization, embedded systems power modeling, and real-time operational benchmarking.

4. ETHICS AND ORIGINALITY DECLARATIONS
This manuscript is an original submission, has not been published previously, and is not under review elsewhere.

Sincerely,

Polisetti Narendra (Corresponding Author)
Department of Computer Science & Engineering
Swarnandhra College of Engineering and Technology (Autonomous)
Seetharampuram, Narasapur-534280, Andhra Pradesh, India
Email: narendraa1996@gmail.com
